% EM structure and teaching revision 20260928
\subsection{分段高斯面与表面跳变}
高斯面必须是闭合曲面；曲面上场强是否相等，要由电荷分布的对称性判断。以半径 $R$、体电荷密度 $\rho$ 的均匀实心球为例，半径 $r$ 的同心球面给出
\[
 E(r)=\begin{cases}\rho r/(3\varepsilon_0),&r<R,\\
 \rho R^3/(3\varepsilon_0r^2),&r>R.
 \end{cases}
\]
这里 $E$ 表示沿径向外的有符号分量。球内包围电荷随 $r^3$ 增加，不能把全体电荷代入球内的高斯定律。由于表面没有额外的电荷薄层，$E$ 在 $r=R$ 连续；其导数可以不连续。
若表面另有总面电荷密度 $\sigma_s$（包含自由与束缚电荷），在界面两侧取薄柱形高斯面，令厚度趋于零，侧面的通量消失，得到
\[
 E_{n,\mathrm{out}}-E_{n,\mathrm{in}}=\frac{\sigma_s}{\varepsilon_0}.
\]
若只给定自由面电荷密度 $\sigma_f$，介质界面的条件应写为 $D_{n,\mathrm{out}}-D_{n,\mathrm{in}}=\sigma_f$，不能把 $\sigma_f$ 直接当作上式的总面电荷。切向分量则由静电场的环量为零而连续。这两条边界关系不要求整个问题球对称，因而也能用于导体表面、带电平面和分层介质的局部检查。
不要把“闭合面净通量为零”误读为“面上处处无场”：外部电荷产生的场可以穿入又穿出同一曲面，净贡献恰好抵消。

